\documentclass[sigplan,review,anonymous=false]{acmart}

%% Rights / copyright — will be set by ACM after acceptance
\setcopyright{none}
\acmConference[Erlang '26]{25th ACM SIGPLAN Erlang Workshop}{August 28, 2026}{Indianapolis, IN, USA}
\acmYear{2026}
\acmDOI{}
\acmISBN{}

%% Packages
\usepackage{lmodern}
\usepackage{listings}
\usepackage{booktabs}

%% Elixir listing style
\lstdefinelanguage{elixir}{
  morekeywords={def,defp,defmodule,do,end,fn,use,case,when,true,false,nil,
    and,or,not,in,forall,property,assert},
  sensitive=true,
  morecomment=[l]{\#},
  morestring=[b]",
  morestring=[b]'
}
\lstset{
  language=elixir,
  basicstyle=\ttfamily\small,
  keywordstyle=\bfseries,
  breaklines=true,
  frame=none,
  xleftmargin=1em,
  aboveskip=0.5em,
  belowskip=0.5em
}

\begin{document}

\title{sim\_ex: Verifying a Concurrent DES Library with PropEr}
\subtitle{Inter-Engine Equivalence, Causal Ordering, and the Bug 700 Sequences Found}

\author{Igor Ostaptchenko}
\email{io@octanix.com}
\affiliation{%
  \institution{Octanix Americas LLC / Wayne State University}
  \city{Detroit}
  \state{MI}
  \country{USA}
}

\begin{abstract}
Discrete event simulation engines on the BEAM have been demonstrated at
scale---Sim-Diasca at EDF, InterSCSimulator for urban traffic---but
their verification has been limited to point tests and trace-based
assertions. We present a three-tier verification methodology for
sim\_ex, a multi-engine DES library for Elixir/OTP: point tests, property
tests with a hybrid stochastic harness, and stateful command-sequence
tests via \texttt{proper\_statem}. Three \texttt{proper\_statem} models
generate approximately 700 adversarial command sequences per run in 1.5
seconds and found three bugs the existing suite had missed, including a
spurious-release protocol error shrunk to four steps. We formulate
inter-engine equivalence as a PropEr property and report results for
strict and statistical equivalence across three of five backends. A
probe against Sim-Diasca's scheduling core (1,600 trials, 8 properties,
no failures) validates transferability. The methodology transfers; the
postconditions are engine-specific.
\end{abstract}

\begin{CCSXML}
<ccs2012>
<concept>
<concept_id>10011007.10011074.10011099.10011102.10011103</concept_id>
<concept_desc>Software and its engineering~Software testing and debugging</concept_desc>
<concept_significance>500</concept_significance>
</concept>
<concept>
<concept_id>10011007.10011074.10011099.10011102.10011104</concept_id>
<concept_desc>Software and its engineering~Formal software verification</concept_desc>
<concept_significance>300</concept_significance>
</concept>
<concept>
<concept_id>10002950.10003714.10003715.10003719</concept_id>
<concept_desc>Mathematics of computing~Stochastic processes</concept_desc>
<concept_significance>100</concept_significance>
</concept>
</ccs2012>
\end{CCSXML}

\ccsdesc[500]{Software and its engineering~Software testing and debugging}
\ccsdesc[300]{Software and its engineering~Formal software verification}
\ccsdesc[100]{Mathematics of computing~Stochastic processes}

\keywords{property-based testing, discrete event simulation, Erlang,
  Elixir, BEAM, PropEr, proper\_statem, inter-engine equivalence}

\maketitle

\section{Introduction}

Discrete event simulation on the BEAM is not new. Sim-Diasca ran
million-actor simulations at EDF in 2010~\cite{boudeville2010};
InterSCSimulator scaled urban traffic models to comparable
sizes~\cite{santana2017}.

What has received less attention is adversarial verification of the
engine itself---the event calendar, the dispatch loop, the resource
protocol, the causal ordering guarantees. We have not found published
work that subjects DES engine internals to property-based or stateful
randomized testing. Point tests catch the bugs the developer
anticipated; the bugs that matter are the ones nobody did.

Property-based testing generates random inputs and checks invariants,
with automatic shrinking to minimal counterexamples~\cite{claessen2000}.
Stateful PBT via \texttt{proper\_statem}~\cite{papadakis2011} generates
random \emph{command sequences}, verifying postconditions after each
step. Applied to Erlang/OTP at Ericsson and Volvo~\cite{arts2006,
hughes2007,hughes2016}, but not, to our knowledge, to a discrete event
simulation engine.

This paper presents the verification methodology for sim\_ex, an
open-source multi-engine DES library for Elixir/OTP with five backends,
a 14-verb DSL, and a Rust NIF engine. Five contributions:

\begin{enumerate}
\item A \textbf{three-tier testing taxonomy}---\emph{clipboard} (point
  tests), \emph{auditor} (property tests with shrinking),
  \emph{saboteur} (stateful command-sequence tests)---and an account of
  how the tiers complement each other.

\item A \textbf{hybrid stochastic harness} that separates deterministic
  invariants (shrinkable) from statistical invariants (not shrinkable) to
  prevent PropCheck's shrinker from chasing Monte Carlo noise.

\item Three \textbf{\texttt{proper\_statem} models} generating
  ${\sim}700$ adversarial command sequences per test run in 1.5~seconds.

\item \textbf{Three bugs found}---a spurious-release protocol error
  (four-step shrunk sequence), a missing pattern clause in the parallel
  engine, and a silent integration failure between the DSL and the Rust
  NIF that returns \texttt{\{:ok, \_\}} with zero events processed.

\item \textbf{Inter-engine equivalence as a PropEr property}, with
  empirical results for strict equivalence (engine vs.\ ETS, 30 trials,
  all identical) and statistical equivalence (engine vs.\ diasca, 60
  trials, within documented tolerance). The equivalence work surfaced an
  undocumented architectural taxonomy: the five backends split into a
  \emph{time group} (float) and a \emph{tick group} (integer with
  $(T, D)$ timestamps).
\end{enumerate}

To test transferability, we extracted the pure scheduling functions from
Sim-Diasca's \texttt{class\_TimeManager.erl} and ran 1,600 PropEr
trials across 8 properties. All passed, as one would expect of code in
current service at EDF since fifteen years.

sim\_ex is open-source under Apache-2.0. \texttt{mix test} reproduces
every result in this paper.

Joe Armstrong was a physicist. Physicists bring \emph{informed
priors}---conservation laws, causality, monotonicity. Armstrong at
Ericsson could not let a phone ring before it was dialed; Boudeville at
EDF whose day job forbids releasing energy before fission. Causality was
not a feature they added; it was the thing they knew how to draw.

But the simulation models that run \emph{on} the BEAM inherit no such
instincts. In a simulation, conservation holds only if someone wrote a
postcondition to check it. Property-based testing is the conservation
law made executable. PropEr enforces what gravity enforces for free. Joe
always advised: solve the hardest problem first. In a simulation engine,
the hardest problem is correctness under inputs no well-formed model
would ever produce---the spurious release, the dangling event, the
calendar that silently reorders---because the developer cannot imagine
what the developer cannot imagine. That is the problem this paper takes
first. The physicist built the machine; we are building the physics.


\section{sim\_ex Architecture}

sim\_ex maps Averill Law's simulation methodology~\cite{law2014} to OTP.
Same model, five backends, no code changes.

\subsection{Entities and the Event Loop}

The fundamental abstraction is the \texttt{Sim.Entity} behaviour:

\begin{lstlisting}
@callback init(config :: map()) :: {:ok, state}
@callback handle_event(event, clock, state)
    :: {:ok, new_state, [new_events]}
@callback statistics(state) :: map()
\end{lstlisting}

An entity receives an event and a clock, updates its state, returns new
events. Calendar: \texttt{:gb\_trees} keyed by \texttt{\{time,
sequence\_number\}} for FIFO tie-breaking. Statistics: Welford's online
algorithm, constant memory.

\subsection{Five Execution Backends}

\textbf{Engine} (default). Tail-recursive loop, single process, Map for
entity states, \texttt{:gb\_trees} calendar. 270K events/sec on PHOLD.

\textbf{ETS}. Same loop, entity states in ETS with write concurrency.
A stepping stone toward parallelism.

\textbf{Diasca}. Implements Sim-Diasca's causal ordering~\cite{boudeville2010}.
Events carry $(T, D)$ timestamps; tick $T{+}1$ begins only at
quiescence. Tags: \texttt{:same\_tick} ${\to}(T,D{+}1)$,
\texttt{:tick} ${\to}(T',0)$, \texttt{:delay} ${\to}(T{+}\delta,0)$.

\textbf{Parallel}. Extends Diasca with concurrent dispatch: events at
$(T, D)$ are causally independent and can be partitioned across a
persistent worker pool.

\textbf{Rust NIF}. Entire simulation in one NIF call: BinaryHeap
calendar, Vec states, \texttt{rand} crate. 6.4M events/sec---$25\times$
the Elixir engine per core. Accepts DSL-compiled step lists only.

\subsection{The Engine Taxonomy}

A discovery from the equivalence work (Section~6): the five engines
split into two groups by time representation.

\begin{table}[h]
\footnotesize
\begin{tabular}{@{}llll@{}}
\toprule
Group & Engines & Clock & Stop \\
\midrule
Time & Eng., ETS & float & stop\_time \\
Tick & Diasca, Par., Rust & \{tick,diasca\} & stop\_tick \\
\bottomrule
\end{tabular}
\end{table}

The protocols are incompatible (Section~5). This taxonomy was not
designed; PropEr surfaced it on the first generated input.

\subsection{The DSL}

A 14-verb language inspired by GPSS~\cite{gordon1961} and
Arena~\cite{kelton2015}, compiled at \texttt{mix compile} time to
\texttt{Sim.Entity} modules:

\begin{lstlisting}
defmodule Barbershop do
  use Sim.DSL
  model :barbershop do
    resource :barber, capacity: 1
    process :customer do
      arrive every: exponential(18.0)
      seize :barber
      hold exponential(16.0)
      release :barber
      depart
    end
  end
end
\end{lstlisting}

\texttt{Sim.Source} detects the clock type by pattern matching and emits
the appropriate event tags. One keyword change switches a DSL model
between \texttt{:engine}, \texttt{:ets}, and \texttt{:diasca}---the
natural input for equivalence testing.

\subsection{Design Principles}

Three principles matter for verification: (1)~\textbf{functional
PRNG}---\texttt{:rand} state threaded through entities; same seed, same
trajectory; CRN for free. (2)~\textbf{Process = entity
(optional)}---the \texttt{Sim.Entity} behaviour is agnostic to hosting
model; same code, five backends. (3)~\textbf{Zero runtime
dependencies}---\texttt{mix test} reproduces every result in this paper.


\section{The Three-Tier Testing Taxonomy}

A simulation engine's inputs are stochastic; its outputs are
distributions. \texttt{grants >= releases} is a conservation law, not a
software requirement. An engine tested without postconditions may be
computationally valid but physically meaningless.

Three strategies: \emph{the clipboard}, \emph{the auditor}, \emph{the
saboteur}.

\subsection{The Clipboard (Point Tests)}

Fixed seed, expected output. Over a hundred tests, under two seconds.
Weakness: the developer must imagine the bug to write the test.

\subsection{The Auditor (Property Tests)}

PropCheck~\cite{alfert2019} generates random M/M/c configurations and
checks invariants: flow conservation, determinism, cross-mode
equivalence, DSL verb semantics. On failure, shrinking isolates the
minimal counterexample. ${\sim}100$ seconds. Weakness: it tests
\emph{scenarios}, not \emph{operation sequences}---it will never release
a job that was never seized.

\subsection{The Saboteur (Stateful Tests)}

\texttt{proper\_statem}~\cite{papadakis2011} generates random command
sequences, verifying postconditions after every step. It releases jobs
never seized, steps engines never initialized. Three models, ${\sim}700$
sequences, 1.5~seconds. The Resource model found the bug the clipboard
and auditor missed (Section~5).

\subsection{The Hybrid Stochastic Harness}

Little's Law holds in expectation, not on finite runs. Written as a
PropCheck property, the shrinker chases noise---reducing run length
until the sample is too small.

Solution: \textbf{deterministic invariants} use PropCheck with
shrinking. \textbf{Stochastic invariants} use a hand-rolled
harness---deterministic seeds, no shrinking, runs long enough for
convergence. Shrink what you can; replay what you can't.


\section{proper\_statem Models}

Three models, different engine layers. Engine state cannot be a PropEr
symbolic variable; the workaround is the process dictionary. Optimistic
\texttt{next\_state} updates (guarded by \texttt{when
is\_number(clock)}) allow meaningful symbolic-phase command generation.

\textbf{Model~1 (Sim.Statham).} Full engine, barbershop scenario.
Commands: \texttt{init}, \texttt{step}, \texttt{run\_n},
\texttt{check}. Postconditions: calendar sortedness, target validity,
flow conservation, clock monotonicity. 200 sequences, 365~ms.

\textbf{Model~2 (Sim.Statham.Resource).} Seize/release protocol in
isolation. The \texttt{release} command is generated only for held job
IDs---the constraint that surfaced the bug. Postconditions: capacity
containment, grant-release balance. 300 sequences, 135~ms.

\textbf{Model~3 (Sim.Statham.Adversarial).} Preemptive resource with
rush/normal order sources. Exercises the generation-counter logic that
invalidates stale hold-complete events. 200 sequences, 150~ms.


\section{Found Bugs}

Three bugs, none affecting a production model, all latent under the
existing test strategy.

\textbf{Bug~1: spurious release.} Four-step shrunk sequence:

{\footnotesize\texttt{init(cap=3,preemptive), seize(1), release(1), release(3)}}

\noindent Job~3 was never seized; the engine accepted the release
silently. \texttt{releases} exceeded \texttt{grants}. Fix: a
\texttt{Map.has\_key?} guard. Adversarial testing defends against bugs
that do not yet exist.

\textbf{Bug~2: parallel engine crash.} Barbershop with
\texttt{mode:~:parallel} crashed in
\texttt{insert\_diasca\_events/5}---the DSL emits float-time
three-tuples; the parallel engine matches only diasca-tagged events.
No test had ever run a DSL model on the parallel engine. A naive bridge
clause caused infinite diasca progression (two BEAM processes at
140\%~CPU, killed with \texttt{SIGKILL}); the fix requires entity-side
tick awareness in the DSL compiler.

\textbf{Bug~3: Rust NIF silent no-op.} \texttt{mode: :rust} returns
\texttt{\{:ok, \%\{events: 0\}\}}. The DSL passes \texttt{:entities};
the Rust engine expects step lists. It silently ignores the unrecognized
options. Worse than a crash: it passes any \texttt{\{:ok, \_\}} smoke
test.

The saboteur does not respect the developer's mental model of what the
system is for. That is the source of its value.


\section{Inter-Engine Equivalence}

Same model, same seed, two engines, different statistics: at least one
is wrong.

\subsection{Equivalence Relation}

\textbf{Strict} (within the time group): bitwise-identical output.

\begin{lstlisting}
r_eng.events == r_ets.events and
  r_eng.stats == r_ets.stats
\end{lstlisting}

\textbf{Statistical} (across the time/tick boundary): integer counts
within max(5, 5\%); float statistics within 35\% relative (rounding
compounds across service stages). Snapshot fields (\texttt{in\_progress},
\texttt{busy}) are excluded---they depend on the exact stop boundary.

\subsection{T1: Within-Time-Group}

Engine vs.\ ETS, random M/M/c scenarios ($\rho \in [0.10, 0.85)$). 30
trials, all identical, 1.5~seconds.

\subsection{T2': Cross-Group DSL Equivalence}

Engine vs.\ ETS vs.\ Diasca, two DSL models (Barbershop and
SplitCombine), 30 seeds each.

\begin{table}[h]
\small
\begin{tabular}{lll}
\toprule
Model & Engine $\leftrightarrow$ ETS & Engine $\leftrightarrow$ Diasca \\
\midrule
Barbershop & All identical & mean\_wait 2--5\% \\
SplitCombine & All identical & mean\_wait up to 25\% \\
\bottomrule
\end{tabular}
\end{table}

180 runs, 3.2~seconds, 0 failures. SplitCombine's divergence:
integer-tick rounding compounding across three stages. A consequence,
not a bug.

\subsection{Discovery: The Engine Taxonomy}

An early equivalence property included \texttt{:parallel}. PropEr's
first input failed: \texttt{:parallel} ignores \texttt{stop\_time} and
operates on \texttt{\{tick, diasca\}}. The property was wrong about what
the engines \emph{are}, not what they \emph{should} do. PBT surfaces
architectural knowledge documented nowhere.


\section{Performance}

All benchmarks: 88-core Xeon E5-2699 v4, OTP~27, Elixir~1.18.3.

\textbf{Event throughput.} Elixir engine: 270K events/sec (PHOLD, 100
LPs). Rust NIF: 6.4M events/sec (barbershop). DSL complexity tax
between simplest and most complex model: $1.6\times$.

\textbf{Parallel replication}---the number that matters:

\begin{table}[h]
\small
\begin{tabular}{lrl}
\toprule
Configuration & Wall time & vs SimPy \\
\midrule
SimPy (sequential) & ${\sim}6{,}300$ ms & $1.0\times$ \\
Elixir parallel (88 cores) & 683 ms & $9.4\times$ \\
Rust NIF parallel (88 cores) & 207 ms & $30\times$ \\
\bottomrule
\end{tabular}
\end{table}

1,000 replications of a 200K-event model. The analysis Law describes as
rarely performed~\cite{law2014} completes in 207~ms. Functional PRNG,
no shared state, no GIL.

\textbf{Verification overhead.} Full suite in ${\sim}110$ seconds. Fast
enough for CI.


\section{Related Work and Conclusion}

\subsection{BEAM-Based Simulation}

Sim-Diasca~\cite{boudeville2010} is the foundational BEAM DES engine.
sim\_ex's contribution is not the engine but the verification.

We extracted two pure functions from Sim-Diasca's
\texttt{class\_TimeManager.erl} (8,289 lines). 1,600 PropEr trials, 8
properties, no failures. Partly explained by architecture: Sim-Diasca's
WOOPER class hierarchy prevents certain bug classes by
construction---actors cannot emit calendar events directly, forge
timestamps, or use unrecognized event tags. We tested the guts; the
architecture protects the surface.

Two design points: Sim-Diasca prevents bugs architecturally; sim\_ex
permits flexibility and relies on testing. The methodology is necessary
for the second, informative for the first.
InterSCSimulator~\cite{santana2017} and Toscano et al.~\cite{toscano2012}
would benefit from the same approach.

\subsection{Property-Based Testing}

Heritage: Claessen and Hughes~\cite{claessen2000}, Papadakis and
Sagonas~\cite{papadakis2011}, Arts et al.~\cite{arts2006},
Hughes~\cite{hughes2007,hughes2016}. Our contribution: the domain
(DES), the hybrid harness, the equivalence framing.

\subsection{Conclusion}

Three tiers, three bugs, one architectural taxonomy the developers did
not know they had. The Sim-Diasca probe passed 1,600 trials, validating
transferability. The postcondition encodes what the developer knows
about the physical system and constrains the simulation to the region
where physics lives. The physicist built the machine; we are building
the physics.

sim\_ex, the test suite, and the Sim-Diasca probe are open-source at
\url{https://github.com/borodark/sim_ex}.


\bibliographystyle{ACM-Reference-Format}
\begin{thebibliography}{10}

\bibitem{boudeville2010}
Olivier Boudeville.
\newblock \emph{Sim-Diasca: Technical Manual of the Sim-Diasca Simulation
  Engine}.
\newblock EDF R\&D, 2010.
\newblock \url{https://github.com/Olivier-Boudeville-EDF/Sim-Diasca}

\bibitem{santana2017}
Eduardo F.~Z. Santana, Nelson Lago, Fabio Kon, and Dejan S. Milojicic.
\newblock InterSCSimulator: Large-scale traffic simulation in smart cities
  using Erlang.
\newblock In \emph{Proc.\ MABS (Multi-Agent Based Simulation)}, LNCS,
  pp.~211--227. Springer, 2018.

\bibitem{claessen2000}
Koen Claessen and John Hughes.
\newblock QuickCheck: A lightweight tool for random testing of Haskell programs.
\newblock In \emph{Proc.\ ICFP}, pp.~268--279. ACM, 2000.

\bibitem{papadakis2011}
Manolis Papadakis and Kostis Sagonas.
\newblock A PropEr integration of types and function specifications with
  property-based testing.
\newblock In \emph{Proc.\ Erlang Workshop}, pp.~39--50. ACM, 2011.

\bibitem{arts2006}
Thomas Arts, John Hughes, Joakim Johansson, and Ulf Wiger.
\newblock Testing telecoms software with Quviq QuickCheck.
\newblock In \emph{Proc.\ Erlang Workshop}, pp.~2--10. ACM, 2006.

\bibitem{hughes2007}
John Hughes.
\newblock QuickCheck testing for fun and profit.
\newblock In \emph{Proc.\ PADL}, LNCS~4354, pp.~1--32. Springer, 2007.

\bibitem{hughes2016}
John Hughes.
\newblock Experiences with QuickCheck: Testing the hard stuff and staying sane.
\newblock In \emph{A List of Successes That Can Change the World}, LNCS~9600,
  pp.~169--186. Springer, 2016.

\bibitem{law2014}
Averill~M. Law.
\newblock \emph{Simulation Modeling and Analysis}.
\newblock McGraw-Hill, 5th edition, 2014.

\bibitem{gordon1961}
Geoffrey Gordon.
\newblock A general purpose systems simulation program.
\newblock In \emph{Proc.\ EJCC}, pp.~87--104. Macmillan, 1961.

\bibitem{kelton2015}
W.~David Kelton, Randall~P. Sadowski, and Nancy~B. Zupick.
\newblock \emph{Simulation with Arena}.
\newblock McGraw-Hill, 6th edition, 2015.

\bibitem{alfert2019}
Klaus Alfert.
\newblock PropCheck: Property-based testing for Elixir.
\newblock \url{https://github.com/alfert/propcheck}, 2019.

\bibitem{toscano2012}
Luca Toscano, Gabriele D'Angelo, and Moreno Marzolla.
\newblock Parallel discrete event simulation with {E}rlang.
\newblock \emph{arXiv:1206.2775}, 2012.

\end{thebibliography}

\end{document}
